\part{Foundations: the bare machine}

\chapter{What a computer really is}

\begin{goals}
  \item What a processor does, and why it only ever does one thing
  \item What memory is, and why it is just a giant list of numbers
  \item What role buses and devices play
  \item Why a powered-on computer with no software does absolutely nothing
\end{goals}

\section{Three parts and a loop}

Strip away the screen, the keyboard, the case and the marketing. What is left is three things:

\begin{enumerate}
  \item A \hi{processor} (CPU) that knows how to execute very simple instructions, very fast.
  \item A \hi{memory} (RAM), which is an extremely long row of numbered boxes, each holding
        one small number.
  \item Some \hi{devices} (disk, keyboard, video card, timer) connected by wires we call
        \emph{buses}.
\end{enumerate}

And a loop. The processor always does the same thing, millions of times per second:

\begin{center}
\begin{tikzpicture}[
  caja/.style={draw=accent,fill=accentlight,rounded corners=3pt,minimum width=3.1cm,
               minimum height=1.05cm,align=center,font=\sffamily\small},
  flecha/.style={-{Stealth[length=3mm]},accent,thick}
]
  \node[caja] (f) {\textbf{1. Fetch}\\{\scriptsize the instruction at EIP}};
  \node[caja,right=1.3cm of f] (d) {\textbf{2. Decode}\\{\scriptsize what am I being asked?}};
  \node[caja,right=1.3cm of d] (e) {\textbf{3. Execute}\\{\scriptsize do it}};
  \draw[flecha] (f) -- (d);
  \draw[flecha] (d) -- (e);
  \draw[flecha] (e.south) -- ++(0,-0.75) -| node[pos=0.25,below,font=\scriptsize\sffamily,text=gray]
        {move on to the next instruction} (f.south);
\end{tikzpicture}
\end{center}

That is all. There is no magic, no intelligence, no ``the computer decides''. There is a register
called \reg{EIP} holding a number: the memory address of the next instruction to execute. The CPU
reads whatever is at that address, interprets it as an order, carries it out, and moves on. If the
order was ``jump to address such-and-such'', \reg{EIP} becomes something else and the loop
continues from there.

\begin{concept}{The only question that matters}
At any moment, the question that explains a computer's behaviour is: \emph{what is in \reg{EIP},
and what is in the memory it points at?} All the work an operating system does amounts, in the
end, to controlling the answer to that question.
\end{concept}

\section{Memory: a list of numbered boxes}

Picture an infinitely long street with houses down one side. Each house has a number (its
\hi{address}) and holds exactly one number between 0 and 255 (a \hi{byte}).

\begin{center}
\begin{tikzpicture}[scale=1,font=\ttfamily\scriptsize]
  \foreach \i/\v in {0/4D,1/5A,2/90,3/00,4/03,5/00,6/00,7/00} {
    \draw[draw=grayborder,fill=graylight] (\i*1.35,0) rectangle (\i*1.35+1.25,0.85);
    \node at (\i*1.35+0.625,0.42) {\v};
    \node[text=gray,font=\ttfamily\tiny] at (\i*1.35+0.625,-0.28) {0x100\i};
  }
  \draw[decorate,decoration={brace,amplitude=4pt},gray]
    (0,1.0) -- (8*1.35-0.1,1.0) node[midway,above,font=\sffamily\scriptsize,text=gray]
    {eight consecutive bytes of memory};
\end{tikzpicture}
\end{center}

That is literally what RAM is. There are no ``variables'', no ``objects'', no ``strings''. There
are boxes with numbers in them. Everything else is an \emph{interpretation} imposed by software.

The same byte \hexa{41} can be:
\begin{itemize}
  \item the number 65,
  \item the letter \cod{'A'} if we read it as ASCII text,
  \item part of the instruction \cod{inc ecx} if the CPU reads it as code,
  \item a dark grey pixel if the graphics card reads it as a colour.
\end{itemize}

\begin{concept}{Data has no type}
A byte in memory does not ``know'' what it is. Type exists only in the programmer's head and in
the code the compiler generates. This apparently philosophical point causes half the faults you
will see in this course: when the system interprets data as code, or writes data over code, there
is nothing stopping it by default.
\end{concept}

\subsection{Why things come in fours}

Our processor is \hi{32-bit}. That means two things at once:

\begin{itemize}
  \item It works comfortably with 32-bit numbers, that is, 4 bytes. We call that group of 4 bytes
        a \term{double word} (\emph{dword}).
  \item Memory addresses are also 32-bit numbers. With 32 bits you can count up to
        $2^{32} = 4{,}294{,}967{,}296$. That is why a 32-bit system can address at most
        \hi{4 GB} of memory. It is not an arbitrary limitation: it is arithmetic.
\end{itemize}

\begin{warn}{Byte order (endianness)}
x86 is \term{little-endian}: storing the number \hexa{12345678} at address 1000 lays the bytes
down backwards from how you wrote them --- \hexa{78} at 1000, \hexa{56} at 1001, \hexa{34} at
1002, \hexa{12} at 1003. It takes getting used to, and it confuses people badly when reading
memory dumps in a debugger. It is not a whim: it simplifies the hardware when you only read the
low byte of a large number.
\end{warn}

\section{Devices and buses}

Processor and memory are not enough: we need to see something on screen and read something from
the keyboard. Devices hang off shared wires called \term{buses}, and the processor talks to them
in two distinct ways.

\subsection{Way 1: memory-mapped}

Some devices appear \emph{as if they were memory}. The video card in text mode is the perfect
example: write a byte to address \hexa{B8000} and it is not stored in RAM --- it travels to the
graphics card and a letter appears in the top-left corner of the screen.

\begin{center}
\begin{tikzpicture}[font=\sffamily\scriptsize,
  b/.style={draw=accent,fill=accentlight,rounded corners=2pt,minimum width=2.4cm,minimum height=0.9cm,align=center}]
  \node[b] (cpu) {CPU};
  \node[b,right=2.6cm of cpu] (ram) {RAM\\{\tiny 0x00000--0x9FFFF}};
  \node[b,below=0.9cm of ram] (vga) {Video card\\{\tiny 0xB8000--0xBFFFF}};
  \draw[-{Stealth[length=2.5mm]},thick,accent] (cpu) -- node[above,text=gray]{\tiny write to 0x1000} (ram);
  \draw[-{Stealth[length=2.5mm]},thick,purple] (cpu.south) |- node[pos=0.75,below,text=gray]{\tiny write to 0xB8000} (vga.west);
\end{tikzpicture}
\end{center}

For the programmer this is convenient: writing to the screen is as simple as assigning through a
pointer. You will see exactly that in Part V.

\subsection{Way 2: I/O ports}

Other devices live in a \emph{separate} address space, the \term{I/O port space}. It holds 65,536
numbered mailboxes reachable only through two special processor instructions: \cod{in} (read from
a port) and \cod{out} (write to a port).

The keyboard, for example, lives at port \hexa{60}. The interrupt controller at \hexa{20} and
\hexa{21}. The disk in the range \hexa{1F0}--\hexa{1F7}.

\filetag{lytlnybl/kernel/interrupts.asm}
\begin{lstlisting}[style=asm]
outb:
    mov al, [esp + 8]   ; the value to send
    mov dx, [esp + 4]   ; the port number
    out dx, al          ; send it
    ret
\end{lstlisting}

Those four lines are the entire mechanism. From C they are wrapped in functions, and from then on
``talking to hardware'' means calling \cod{outb(port, value)}.

\begin{analogy}
Memory-mapped I/O is like leaving a note on someone's desk: you write it, and the note is there.
Ports are like an intercom: you have to dial the office number (\reg{DX}), say the phrase
(\reg{AL}) and press the button (\cod{out}). The effect is not stored anywhere; it simply
\emph{happens}.
\end{analogy}

\section{Why a powered-on computer does nothing}

Here is the question almost nobody asks, and which is the starting point of this whole course:
when you press the power button, what does the processor execute?

The CPU wakes up with \reg{EIP} pointing at a fixed address, hard-wired into the silicon. There is
no RAM at that address (RAM has just powered on and contains garbage): there is a read-only memory
chip holding the \term{firmware}, which on a classic PC is called the \hi{BIOS}.

In other words: the computer does not ``boot itself''. It boots because the manufacturer put a
program in a chip, at exactly the address the processor looks at when it wakes. And that program,
after checking the hardware, has to find \emph{another} program on disk and hand over to it. And
that other program has to load the operating system.

\begin{concept}{The boot chain}
Power on $\rightarrow$ firmware (BIOS) $\rightarrow$ boot sector (512 bytes on disk)
$\rightarrow$ operating system kernel $\rightarrow$ user programs. Each link has exactly one job:
leave the next one in a position to start. All of Part III is this chain.
\end{concept}

\exercise{Without reading ahead, reason it out: if the boot sector is exactly 512 bytes, why can
it not contain the whole operating system? What is the minimum it would have to do?}

% ============================================================
\chapter{Binary, hexadecimal and bits}

\begin{goals}
  \item Read and write hexadecimal fluently
  \item Understand the bit operations: AND, OR, XOR, NOT and shifts
  \item Recognise the \cod{| FLAG}, \cod{\& \textasciitilde FLAG} and \cod{\& MASK} idioms that appear all over the code
\end{goals}

This chapter is short and purely practical. If you are already comfortable with bits, skim it and
move on. If not, read it slowly: without this, 60\,\% of operating system code looks like noise.

\section{Why hexadecimal}

A byte is 8 bits. In binary a byte reads \cod{11010110}: exact but unreadable. In decimal, 214:
readable but telling you nothing about the bits.

Hexadecimal (base 16) is the perfect middle ground because \hi{each hex digit corresponds to
exactly 4 bits}. A byte is always two hex digits. No more, no less.

\begin{center}
\begin{tabularx}{0.92\textwidth}{@{}llllX@{}}
\toprule
\textbf{Hex} & \textbf{Binary} & \textbf{Decimal} & \textbf{Hex} & \textbf{Binary / Decimal} \\
\midrule
\cod{0} & \cod{0000} & 0 & \cod{8} & \cod{1000} \quad 8 \\
\cod{1} & \cod{0001} & 1 & \cod{9} & \cod{1001} \quad 9 \\
\cod{2} & \cod{0010} & 2 & \cod{A} & \cod{1010} \quad 10 \\
\cod{3} & \cod{0011} & 3 & \cod{B} & \cod{1011} \quad 11 \\
\cod{4} & \cod{0100} & 4 & \cod{C} & \cod{1100} \quad 12 \\
\cod{5} & \cod{0101} & 5 & \cod{D} & \cod{1101} \quad 13 \\
\cod{6} & \cod{0110} & 6 & \cod{E} & \cod{1110} \quad 14 \\
\cod{7} & \cod{0111} & 7 & \cod{F} & \cod{1111} \quad 15 \\
\bottomrule
\end{tabularx}
\end{center}

With that table, \hexa{D6} translates at a glance: \cod{D} = \cod{1101}, \cod{6} = \cod{0110},
so \cod{11010110}. And back again just as fast.

\begin{concept}{Round numbers in hexadecimal}
\hexa{1000} = 4096 = 4 KB. And that is no coincidence: \hi{4096 bytes is the size of a
memory page} on x86. That is why you see addresses like \hexa{B8000}, \hexa{100000} or
\hexa{C0000000} everywhere: they end in three hex zeros, which means they are page-aligned. When
you see an address ending in \cod{000}, think ``start of a page''.
\end{concept}

Some figures worth memorising, because they appear constantly:

\begin{center}
\begin{tabularx}{0.92\textwidth}{@{}l l X@{}}
\toprule
\textbf{Value} & \textbf{Decimal} & \textbf{What it is in this course} \\
\midrule
\hexa{1000}      & 4,096         & one page / one memory frame \\
\hexa{200}       & 512           & one disk sector, and the filesystem block size \\
\hexa{B8000}     & 753,664       & start of video memory in text mode \\
\hexa{7C00}      & 31,744        & where the BIOS loads the boot sector \\
\hexa{100000}    & 1,048,576     & 1 MB; end of ``low memory'' \\
\hexa{400000}    & 4,194,304     & 4 MB; where user code starts \\
\hexa{C0000000}  & 3,221,225,472 & 3 GB; where the kernel heap starts \\
\hexa{FFFFFFFF}  & 4,294,967,295 & the highest possible 32-bit address \\
\bottomrule
\end{tabularx}
\end{center}

\section{The five bit operations}

\subsection{AND (\cod{\&}): switching off and asking}

\cod{a \& b} leaves a 1 only where \emph{both} had a 1. It has two uses:

\begin{lstlisting}[style=c]
/* 1. ASK whether a bit is set */
if (status & ATA_STATUS_BSY) { /* the disk is busy */ }

/* 2. KEEP part of a number (masking) */
uint32_t base_address = entry & 0xFFFFF000;     /* the top 20 bits */
uint16_t table_index  = (virtual >> 12) & 0x3FF; /* 10 bits */
\end{lstlisting}

The second form is everywhere in paging. \hexa{FFFFF000} has the top 20 bits set and the low 12
clear: it keeps a page address and throws away the permission bits. \hexa{3FF} is ten ones: it
keeps an index from 0 to 1023.

\subsection{OR (\cod{|}): switching on}

\cod{a | b} leaves a 1 where \emph{either} had one. This is how a value is built from flags:

\begin{lstlisting}[style=c]
map_page(dir, virtual, physical, PAGE_PRESENT | PAGE_WRITABLE | PAGE_USER);
\end{lstlisting}

Each constant has exactly one bit set, so \cod{|} combines them without collision. This idiom ---
define flags as powers of two, combine with \cod{|} --- appears dozens of times.

\subsection{NOT (\cod{\textasciitilde}) with AND: switching one bit off}

\begin{lstlisting}[style=c]
bitmap[i / 32] &= ~(1 << (i % 32));   /* clear bit i */
(*table)[index] &= ~(PAGE_PRESENT);   /* mark the page as not present */
\end{lstlisting}

\cod{\textasciitilde X} inverts every bit; ANDing with that clears exactly the bits that were set
in \cod{X} and leaves everything else alone. It is the canonical way to say ``remove this flag and
touch nothing else''.

\subsection{Shifts (\cod{<<} and \cod{>>})}

Shifting left by $n$ multiplies by $2^n$; shifting right divides. But their most frequent use is
\emph{placing} and \emph{extracting} fields:

\begin{lstlisting}[style=c]
/* Placing: build a VGA attribute byte from two 4-bit colours */
terminal->color = ((uint8_t)background << 4) | (uint8_t)foreground;

/* Extracting: split a virtual address into its three fields */
uint16_t directory_index = virtual >> 22;            /* bits 31..22 */
uint16_t table_index     = (virtual >> 12) & 0x3FF;  /* bits 21..12 */
uint16_t offset          = virtual & 0xFFF;          /* bits 11..0  */
\end{lstlisting}

\begin{concept}{The language of bits}
Five idioms cover almost everything: \cod{x | F} sets, \cod{x \&= \textasciitilde F} clears,
\cod{x \& F} asks, \cod{x \& MASK} trims, \cod{x >> n} extracts. Recognise them on sight and
systems code stops looking cryptic overnight.
\end{concept}

\exercise{The physical memory manager keeps one bit per 4 KB page. The function that marks page
$i$ as used is \cod{bitmap[i / 32] |= (1 << (i \% 32));}. Explain why it divides by 32, and why
it takes the remainder modulo 32.}

% ============================================================
\chapter{Inside the x86 processor}

\begin{goals}
  \item What registers are and what each one is for
  \item What flags are and why conditional jumps exist
  \item What privilege rings are (ring 0 and ring 3)
  \item What the control registers \reg{CR0}, \reg{CR2} and \reg{CR3} do
\end{goals}

\section{Registers: the processor's own variables}

Memory is huge but slow. Inside the chip itself there are a handful of extremely fast boxes: the
\hi{registers}. On 32-bit x86 there are eight general-purpose ones, each 32 bits wide.

\begin{center}
\begin{tabularx}{\textwidth}{@{}l l X@{}}
\toprule
\textbf{Register} & \textbf{Name} & \textbf{Typical use} \\
\midrule
\reg{EAX} & Accumulator & Operation results; function return value; \textbf{system call number} \\
\reg{EBX} & Base        & General purpose; first argument in our system calls \\
\reg{ECX} & Counter     & Loop counter; second argument in our system calls \\
\reg{EDX} & Data        & I/O operations (holds the port); third argument \\
\reg{ESI} & Source      & Source pointer in memory copies \\
\reg{EDI} & Destination & Destination pointer in memory copies \\
\reg{EBP} & Base pointer& Marks the start of the current function's frame \\
\reg{ESP} & Stack pointer & Points at the top of the stack. \textbf{Special: push and pop modify it} \\
\bottomrule
\end{tabularx}
\end{center}

The names come from historical convention; apart from \reg{ESP} (and partly \reg{EBP}), the
processor lets you use them for whatever you like.

\subsection{One register, three names}

For compatibility with processors from the 1970s and 80s, each register can be used whole or in
pieces:

\begin{center}
\begin{tikzpicture}[font=\ttfamily\small]
  \draw[draw=grayborder,fill=accentlight] (0,0) rectangle (8,0.9);
  \node at (4,0.45) {EAX \ (32 bits)};
  \draw[draw=grayborder,fill=purplelight] (4,-1.1) rectangle (8,-0.2);
  \node at (6,-0.65) {AX \ (16 bits)};
  \draw[draw=grayborder,fill=greenlight] (4,-2.2) rectangle (6,-1.3);
  \node at (5,-1.75) {AH};
  \draw[draw=grayborder,fill=amberlight] (6,-2.2) rectangle (8,-1.3);
  \node at (7,-1.75) {AL};
  \node[font=\sffamily\scriptsize,text=gray,anchor=west] at (8.3,0.45) {bits 31--0};
  \node[font=\sffamily\scriptsize,text=gray,anchor=west] at (8.3,-0.65) {bits 15--0};
  \node[font=\sffamily\scriptsize,text=gray,anchor=west] at (8.3,-1.75) {bits 15--8 and 7--0};
\end{tikzpicture}
\end{center}

This explains code that is otherwise baffling. In the bootloader you will see:

\begin{lstlisting}[style=asm]
mov ah, 0Eh    ; BIOS function: print character
mov al, 'H'    ; the character to print
int 10h        ; call the BIOS video service
\end{lstlisting}

Two different values are being loaded into \emph{two halves of the same register} \reg{AX},
because that is the interface the BIOS expects.

\section{Flags: the invisible result}

Every arithmetic or logic operation, besides its result, updates a special register called
\reg{EFLAGS}. Each bit answers a question about what just happened:

\begin{center}
\begin{tabularx}{0.95\textwidth}{@{}l l X@{}}
\toprule
\textbf{Flag} & \textbf{Bit} & \textbf{Means} \\
\midrule
\reg{ZF} (Zero)      & 6  & The result was exactly zero \\
\reg{CF} (Carry)     & 0  & There was a carry, or an error; the BIOS uses it to signal failure \\
\reg{SF} (Sign)      & 7  & The result is negative \\
\reg{IF} (Interrupt) & 9  & \textbf{Interrupts are enabled} \\
\bottomrule
\end{tabularx}
\end{center}

A conditional jump does not look at the result: it looks at the flags. \cod{cmp a, b} internally
performs a subtraction and throws the result away, keeping only the flags; \cod{je} (``jump if
equal'') really jumps if \reg{ZF} is set.

\begin{warn}{\reg{IF} and the value \hexa{202}}
All over the process code you will see \cod{eflags = 0x202}. In binary that is
\cod{0000 0010 0000 0010}: bit 1 is always set (reserved, the processor requires it) and bit 9 is
\reg{IF}. So \hexa{202} means ``a clean state \hi{with interrupts enabled}''. If a process
started with \hexa{002}, it would never yield control back to the scheduler and would hang the
system.
\end{warn}

\section{Rings: why not everyone may do everything}

The x86 processor has four privilege levels, called \term{rings}, numbered 0 to 3. In practice
almost every system uses only two:

\begin{center}
\begin{tikzpicture}[font=\sffamily\small]
  \fill[accentlight,draw=accent] (0,0) circle (2.6);
  \fill[redlight,draw=red] (0,0) circle (1.3);
  \node[text=red,font=\sffamily\bfseries\small] at (0,0.15) {Ring 0};
  \node[text=red,font=\sffamily\scriptsize] at (0,-0.35) {the kernel};
  \node[text=accent,font=\sffamily\bfseries\small] at (0,2.0) {Ring 3};
  \node[text=accent,font=\sffamily\scriptsize] at (0,1.6) {the programs};

  \node[align=left,anchor=west,font=\scriptsize,text=gray] at (3.2,1.2)
    {\textbf{\color{red}Ring 0 may:}\\ run \texttt{cli}, \texttt{lgdt}, \texttt{in}, \texttt{out}\\
     change \texttt{CR0}/\texttt{CR3}\\ access any page};
  \node[align=left,anchor=west,font=\scriptsize,text=gray] at (3.2,-1.1)
    {\textbf{\color{accent}Ring 3 may:}\\ arithmetic, jumps, calls\\ read/write \emph{its own} memory\\
     and little else};
\end{tikzpicture}
\end{center}

If a ring 3 program tries to execute \cod{cli} (disable interrupts), the processor does not run
it: it raises a \term{general protection fault}, and the kernel decides what to do with that
program. Usually, kill it.

\begin{concept}{Isolation is not done in software}
This is the key idea of the whole course. An operating system does not stop a program from doing
damage by ``watching it'': it \hi{configures the hardware so the program is physically
incapable of it}. The kernel sets the rules once (in the GDT, in the page tables) and from then on
the silicon enforces them, at zero cost. If the kernel sets those rules wrongly, the isolation
simply does not exist.
\end{concept}

\section{The control registers}

Besides the general-purpose registers, some registers control how the processor operates. This
course uses three:

\begin{description}
  \item[CR0] Master switches. Bit 0 (\emph{PE}) turns on protected mode. Bit 31 (\emph{PG}) turns
    on paging. Setting those two bits, in the right order, is literally the moment the system
    becomes a modern one.
  \item[CR2] When a page fault occurs, the processor leaves \emph{the address that was accessed}
    here. It is the first question any fault handler asks.
  \item[CR3] Holds the physical address of the active page directory. \hi{Changing
    \reg{CR3} means changing address space}, which makes it the heart of a context switch.
\end{description}

\filetag{lytlnybl/memory/vmm.asm}
\begin{lstlisting}[style=asm]
set_cr3:
    mov eax, [esp + 4]
    mov cr3, eax        ; switch the active address space
    ret

set_cr0:
    mov eax, cr0
    or eax, 0x80000000  ; bit 31 = PG: enable paging
    mov cr0, eax
    ret
\end{lstlisting}

Three instructions. With them, the processor stops seeing physical memory directly and starts
seeing a translated world. Part VIII explains what that means.

\exercise{Why do you think \reg{CR3} holds a \emph{physical} address rather than a virtual one?
Think about what would happen if the processor had to translate the address of the page directory
using the page directory itself.}

% ============================================================
\chapter{The memory of a PC}

\begin{goals}
  \item How the stack works and why it grows downwards
  \item What a stack frame is, and exactly what happens on a function call
  \item How a PC's memory is laid out at boot
  \item What a stack overflow is and why it is so destructive
\end{goals}

\section{The stack}

The \term{stack} is a region of memory with a very simple discipline: the last thing you put in is
the first thing you take out. The processor supports it directly with two instructions:

\begin{itemize}
  \item \cod{push value}: subtracts 4 from \reg{ESP} and writes the value there.
  \item \cod{pop dest}: reads the value at \reg{ESP} and adds 4 to \reg{ESP}.
\end{itemize}

Notice the crucial detail: \hi{\cod{push} subtracts}. The stack grows towards \emph{lower}
addresses. It is a historical convention with an enormous practical consequence: if you place the
stack just above your data, growing it overwrites them.

\begin{center}
\begin{tikzpicture}[font=\ttfamily\scriptsize]
  \draw[draw=grayborder] (0,0) rectangle (3.4,4.6);
  \node[font=\sffamily\scriptsize,text=gray] at (1.7,4.85) {high addresses};
  \node[font=\sffamily\scriptsize,text=gray] at (1.7,-0.3) {low addresses};

  \fill[graylight] (0,3.6) rectangle (3.4,4.6);
  \node[font=\sffamily\scriptsize,text=gray] at (1.7,4.1) {(unused)};

  \fill[accentlight] (0,2.6) rectangle (3.4,3.6);
  \node at (1.7,3.1) {0x00000042};
  \fill[accentlight] (0,1.6) rectangle (3.4,2.6);
  \node at (1.7,2.1) {0x0000ABCD};
  \fill[accentlight] (0,0.6) rectangle (3.4,1.6);
  \node at (1.7,1.1) {0x00009000};

  \draw[draw=grayborder] (0,0.6) -- (3.4,0.6);
  \node[font=\sffamily\scriptsize,text=gray] at (1.7,0.3) {free space};

  \draw[-{Stealth[length=2.5mm]},thick,red] (4.6,1.1) -- (3.5,1.1);
  \node[anchor=west,font=\sffamily\scriptsize,text=red] at (4.7,1.1) {ESP};
  \draw[-{Stealth[length=2.5mm]},thick,accent] (1.7,-0.75) -- (1.7,0.35);
  \node[anchor=north,font=\sffamily\scriptsize,text=accent] at (1.7,-0.8) {grows this way};
\end{tikzpicture}
\end{center}

\subsection{What happens on a function call}

When code executes \cod{call my\_function}, this happens:

\begin{enumerate}
  \item The arguments are pushed onto the stack (on 32-bit x86, right to left).
  \item \cod{call} pushes the \hi{return address}: the address of the following instruction.
        Then it jumps.
  \item The function pushes \reg{EBP} and does \cod{mov ebp, esp}, so \reg{EBP} stays fixed at the
        start of its frame while \reg{ESP} keeps moving.
  \item The function subtracts from \reg{ESP} to reserve its local variables.
  \item On return, \cod{leave} undoes the above and \cod{ret} pops the return address and jumps to
        it.
\end{enumerate}

\begin{lstlisting}[style=asm]
vga_text_writeline:
    push ebp            ; save the previous frame
    mov  ebp, esp       ; establish the new frame
    sub  esp, 8         ; reserve 8 bytes of locals
    ; ... body of the function ...
    leave               ; mov esp,ebp ; pop ebp
    ret                 ; jump to the return address
\end{lstlisting}

\begin{danger}{Why a stack overflow is catastrophic}
The return address lives \emph{on the stack}, among the data. If a function writes past the end of
a local array, it clobbers that address. When the \cod{ret} arrives, the processor jumps to
whatever that corrupted value says: usually, nowhere.

In a normal program that produces a segmentation fault and the operating system kills it.
\hi{In a kernel there is nobody above to do the killing.} The processor tries to execute
garbage, faults, tries to handle the fault, faults again, and on the third failure it resets
itself. Part XIII walks through a real case of exactly this, step by step.
\end{danger}

\section{The memory map of a PC}

When the system boots, memory is neither empty nor uniform. There are regions reserved for
historical reasons going back to 1981:

\begin{center}
\begin{tikzpicture}[font=\sffamily\scriptsize,x=1cm,y=1cm]
  \def\w{7.2}
  \newcommand{\banda}[5]{%
    \fill[#4,draw=grayborder] (0,#1) rectangle (\w,#2);
    \node[align=center,font=\sffamily\scriptsize] at (\w/2,{(#1+#2)/2}) {#3};
    \node[anchor=west,font=\ttfamily\tiny,text=gray] at (\w+0.15,#1) {#5};
  }
  \banda{0}{0.7}{Interrupt vector table and BIOS data}{amberlight}{0x00000}
  \banda{0.7}{1.5}{\textbf{Free memory} \\ {\tiny the kernel is loaded here (0x9000)}}{greenlight}{0x00500}
  \banda{1.5}{2.1}{Boot sector area {\tiny (0x7C00)}}{purplelight}{0x07C00}
  \banda{2.1}{3.0}{\textbf{Free memory} \\ {\tiny E820 map at 0x5000, bitmap at 0x100000}}{greenlight}{0x07E00}
  \banda{3.0}{3.7}{Video memory {\tiny (text at 0xB8000)}}{accentlight}{0xA0000}
  \banda{3.7}{4.4}{BIOS and option ROMs}{graylight}{0xC0000}
  \banda{4.4}{5.6}{\textbf{Extended memory} \\ {\tiny all modern RAM lives here}}{greenlight}{0x100000}
  \node[anchor=west,font=\ttfamily\tiny,text=gray] at (\w+0.15,5.6) {0xFFFFFFFF};
\end{tikzpicture}
\end{center}

The holes in that list matter: if the memory manager hands out a page that is really video memory,
the system starts writing data and strange symbols appear on screen. Worse: if it hands out a page
holding its own code, it overwrites itself.

\begin{concept}{This is why E820 exists}
Since the real map depends on the machine, the BIOS offers a service (function \hexa{E820}) that
returns a list of regions: address, size and type (1 = usable, anything else = do not touch). The
lytlnyblOS bootloader asks for that list and stores it at \hexa{5000} so the kernel can read it
later. You will see that code in Chapter 12.
\end{concept}

\exercise{Looking at the map, why do you think the boot sector is loaded at \hexa{7C00} rather
than at \hexa{0}? Hint: look at what lives at \hexa{0}, and how much room is left above
\hexa{7C00} before the next reserved block.}

% ============================================================
\chapter{Talking to hardware}

\begin{goals}
  \item How \cod{in} and \cod{out} are implemented and called from C
  \item Why you sometimes have to wait for hardware
  \item What polling is and why it is a bad long-term idea
\end{goals}

\section{The port functions}

All conversation with devices goes through four tiny assembly functions, because C has no way to
express the \cod{in} and \cod{out} instructions.

\filetag{lytlnybl/kernel/interrupts.asm}
\begin{lstlisting}[style=asm]
outb:                       ; void outb(uint16_t port, uint8_t value)
    mov al, [esp + 8]
    mov dx, [esp + 4]
    out dx, al
    ret

inb:                        ; uint8_t inb(uint16_t port)
    mov dx, [esp + 4]
    in  al, dx              ; the result lands in AL, and AL is part of EAX,
    ret                     ; which is where C expects a return value
\end{lstlisting}

The \cod{inb} detail is worth pausing on: there is no instruction that ``returns'' a value. The C
calling convention simply says the return value travels in \reg{EAX}, and \cod{in al, dx} already
leaves the byte in \reg{AL}, the low byte of \reg{EAX}. The function works by construction.

The 16-bit versions (\cod{inw}/\cod{outw}) are identical with \reg{AL} replaced by \reg{AX}, and
are used for the disk, which transfers two bytes at a time.

\section{Waiting for hardware}

Devices are vastly slower than the processor. Send two commands back to back to an old chip and
the second may arrive before it has digested the first. There are two ways to deal with that.

\subsection{A blind pause}

\begin{lstlisting}[style=asm]
io_wait:
    out 0x80, al    ; write to a port that does nothing
    ret             ; the bus round trip burns the time we need
\end{lstlisting}

Port \hexa{80} is connected to nothing useful on a PC. Writing to it has no effect, but the bus
cycle takes long enough (a few microseconds) to give a slow chip time to catch up. It is an old,
ugly, perfectly effective trick. It is used when reprogramming the interrupt controller.

\subsection{Asking until it is ready}

The disk is more sophisticated: it has a status register you can poll.

\filetag{lytlnybl/kernel/drivers/ata.c}
\begin{lstlisting}[style=c]
static void ata_wait_bsy() {
  uint8_t status;
  do {
    status = inb(ATA_STATUS);
  } while (status & ATA_STATUS_BSY);   /* spin until it stops being busy */
}
\end{lstlisting}

This is \term{polling}: asking in a loop. It is simple and correct, but while the loop spins
\hi{the processor does nothing else}. During a disk read the whole system is frozen.

\begin{warn}{Polling versus interrupts}
The alternative is to ask the device to \emph{signal} when it finishes, via an interrupt, and run
another process meanwhile. That is what real systems do, and it is considerably more complex.
lytlnyblOS uses interrupts for keyboard and timer but polling for disk: a very typical mix in
small systems. It is a conscious trade-off, not an oversight.
\end{warn}

\begin{summary}
You now have the vocabulary of the machine: registers, flags, stack, addresses, ports and rings.
Part II steps up a level and answers the underlying question: given that hardware, what exactly
must an operating system contribute, and why?
\end{summary}
